\documentclass[11pt]{article}
\usepackage[T1]{fontenc}
\usepackage[margin=1in]{geometry}
\usepackage{booktabs}
\usepackage{muhib-gpt2}
\title{Muhib's GPT-2, Compiled}
\author{Muhib Waqar\\University of Waterloo}
\date{}
\begin{document}
\maketitle

\noindent This page was not written in advance. While \texttt{pdflatex} compiled it,
TeX ran the full 124M-parameter GPT-2 forward pass (all 12 layers and the
50{,}257-way output head) and picked each next token itself.\footnote{All of the
mathematics is eTeX \texttt{\string\numexpr} integer arithmetic: about 124 million
multiply--adds per generated token, plus layer norm, softmax, GELU and the BPE
tokenizer, all written as TeX macros. Shell escape for this run:
\ifnum\pdfshellescape=0 off\else on\fi. No Lua.}

\bigskip
\begin{quote}\Large
\muhib[12]{Hello, my name is Muhib}
\end{quote}
\bigskip

\noindent Projected compile time, according to the author: instant. Projected
date this TeX file reaches AGI: next Tuesday.\par

\bigskip
\noindent The prompt was split into \gptnumprompt\ tokens by TeX's own BPE:
\gptpromptids.

\newcount\row \def\rows{}
\newcommand\gptrow[1]{#1 & \gptgenid{#1} & \texttt{"\gptgentext{#1}"} & \gptgenseconds{#1}\\}
\loop\ifnum\row<\gptnumgenerated\relax \advance\row 1
  \edef\rows{\unexpanded\expandafter{\rows}\noexpand\gptrow{\the\row}}\repeat
\begin{center}
\begin{tabular}{rrlr}
\toprule
\# & token id & text & seconds since start \\
\midrule
\rows
\bottomrule
\end{tabular}
\end{center}
\end{document}
